\section{Projection fields and transfer between endpoints}
\label{sec:fields}

We now turn the parameters of Proposition~\ref{prop:parameters} into a
weighted shift.  The purpose of this section is to convert invariance under
that shift into an identity between projection fields at two different
points of the circle.  We first obtain the identity along positive rotation
orbits, and then use the sampling estimates to reach every translation by
an integer multiple of $1/q_A$, for some $A\ge1$.

\subsection{Fields determined by scalar multipliers}

We first isolate the elementary field representation used here and in the
direct weighted-shift construction.  It depends only on commutation with
scalar multipliers.

\begin{lemma}\label{lem:scalar-multiplier-fields}
Let $(Y,\nu)$ be a probability space, let $K=\ell^2(\mathbb Z)$, and let
$A\in B(L^2(Y,\nu;K))$ commute with every scalar multiplication operator
$M_g\zeta(y)=g(y)\zeta(y)$, $g\in L^\infty(Y,\nu)$.
There is a field $A(y)\in B(K)$, unique
almost everywhere, such that
\[
 (A\zeta)(y)=A(y)\zeta(y),\qquad
 \|A(y)\|\le\|A\|\quad\text{almost everywhere}.
\]
The field is measurable in the sense that $y\mapsto A(y)v$ is measurable
for every $v\in K$.  Products and adjoints of such operators are computed
fiberwise.  In particular, an orthogonal projection is represented by
orthogonal projections outside one null set.
\end{lemma}

\begin{proof}
If $B\in B(L^2(Y,\nu))$ commutes with all multipliers, put $b=B1$.
For every measurable set $E$,
\[
 B\mathbf1_E=\mathbf1_E b,\qquad
 \int_E|b|^2\,d\nu\le\|B\|^2\nu(E).
\]
Thus $|b|\le\|B\|$ almost everywhere.  The identity $Bf=bf$ follows
first for simple functions and then, by density, for all $f\in L^2(Y)$.
Apply this observation to the matrix entries of $A$ between the coordinate
copies of $L^2(Y)$, obtaining measurable functions $a_{ij}$.

Let $v$ range over the countable set of finitely supported vectors with
rational real and imaginary parts.  Commutation with $M_{\mathbf1_E}$ and
the operator bound give
\[
 \int_E\sum_{i\in\mathbb Z}
       \left|\sum_j a_{ij}(y)v_j\right|^2d\nu(y)
 \le\|A\|^2\|v\|^2\nu(E).
\]
Consequently the integrand is at most $\|A\|^2\|v\|^2$ almost everywhere.
Outside the union of these countably many null sets, the matrix defines
an operator $A(y)$ on $K$ with norm at most $\|A\|$.  Its columns, and
then its values on every vector, are measurable by finite-coordinate
approximation.  The represented field agrees with $A$ on finite-coordinate
simple functions, hence on all of $L^2(Y;K)$.

Equality of two such fields follows by applying their represented
operators to each constant basis field.  Products act fiberwise by the
representation formula, and the entries of the adjoint field are
$\overline{a_{ji}}$.  This proves the assertions about products and
adjoints.  The identities $p=p^*=p^2$ then give orthogonal projections
on a common conull set.
\end{proof}

\subsection{The tracial representation}

We use the standard crossed-product representation for a
probability-preserving action; see
\citet[Section~1.4 and Proposition~1.4.5]{AnantharamanPopa}.
We verify below the field and trace properties needed in the proof.

Write $X=\mathbb R/\mathbb Z$, let $\mu$ be Haar probability measure, and put
$\sigma x=x+\theta$.  On
$\mathcal H=L^2(X,\mu;\ell^2(\mathbb Z))$, let $S e_k=e_{k+1}$ and define
\[
 U(x)=S,\qquad
 \pi(a)(x)e_k=a(\sigma^k x)e_k\quad(a\in L^\infty(X,\mu)).
\]
For $V=\pi(x\mapsto e^{2\pi i x})$, these operators satisfy
$VU=e^{2\pi i\theta}UV$.
We also write $e_0$ for the constant field with this value.
The vector $U^mV^n e_0$ is the field $x\mapsto e^{2\pi i nx}e_m$.
These vectors span a dense subspace by ordinary Fourier theory, and
$\langle U^mV^n e_0,e_0\rangle$ equals $1$ for $(m,n)=(0,0)$ and $0$
otherwise.  The commutation relation reduces every product of monomials
to a scalar multiple of a monomial, so these moments identify all inner
products of polynomial vectors with the prescribed tracial cyclic
representation.  The resulting unitary identifies the generated algebra
$M$ with $R_\theta$ and its faithful normal trace with $\tau$.
Measurable functional calculus for $V$ places every $\pi(a)$ in $M$.

Every generator commutes with scalar multiplication by a function of $x$.
Lemma~\ref{lem:scalar-multiplier-fields} therefore represents every $A\in M$
by an essentially bounded measurable field.  The generators also commute
with the unitary $J$ given by $(J\zeta)(x)=S\zeta(\sigma x)$:
for $\pi(a)$ the coordinate calculation is
$a(\sigma^{k-1}(\sigma x))=a(\sigma^kx)$.
These commutations pass to their von Neumann algebra.  Testing $AJ=JA$ on
the constant basis fields gives
\begin{equation}\label{eq:field-covariance}
 A(\sigma x)=S^{-1}A(x)S\quad\text{for almost every }x.
\end{equation}
The vector formula for the trace is now
\begin{equation}\label{eq:field-trace}
 \tau(A)=\int_X\langle A(x)e_0,e_0\rangle\,d\mu(x)
 \qquad(A\in M).
\end{equation}
We choose Borel versions for projection fields and discard null sets and
their countably many rotation translates when necessary.  Irrational
rotation is ergodic: all Fourier coefficients of nonzero index vanish
for an invariant $L^2$ function.  We use equidistribution
in the form established at the beginning of Section~\ref{sec:parameters}.

The following consequence of covariance will also apply to the product
model.  It ensures that a nontrivial invariant projection has a nonzero
zeroth coordinate at almost every endpoint, as does its complement.

\begin{lemma}\label{lem:positive-diagonal}
Let $\sigma$ be an invertible ergodic probability-preserving transformation
of $(Y,\nu)$.  On $L^2(Y,\nu;\ell^2(\mathbb Z))$, define
$(U\zeta)(y)=S\zeta(y)$ and
$\pi(a)(y)e_k=a(\sigma^ky)e_k$, where $Se_k=e_{k+1}$.
Suppose a von Neumann algebra $M$ containing $U$ and $\pi(L^\infty(Y))$
consists of measurable fields with covariance
$A(\sigma y)=S^{-1}A(y)S$ and has faithful normal tracial state
\[
 \tau(A)=\int_Y\langle A(y)e_0,e_0\rangle\,d\nu(y).
\]
Let $f\in L^\infty(Y,\nu)$ be positive almost everywhere, put $T=U\pi(f)$,
and let $p\in M$ be an orthogonal projection with $(1-p)Tp=0$.
If $0<\tau(p)<1$, then $0<\langle p(y)e_0,e_0\rangle<1$ almost everywhere.
\end{lemma}

\begin{proof}
Put $r(y)=\langle p(y)e_0,e_0\rangle$ and
$D_y=\operatorname{diag}_k f(\sigma^ky)$.
Invariance and covariance give
$D_y\operatorname{ran}p(y)\subseteq\operatorname{ran}p(\sigma y)$.
If $r(\sigma y)=0$, every vector in the latter range has zero zeroth
coordinate.  Since $(D_yv)_0=f(y)v_0$ and $f(y)>0$, every vector in the
former range also has zero zeroth coordinate, so $r(y)=0$.
If $r(y)=1$, then $e_0$ belongs to the former range, and its image
$f(y)e_0$ shows that $r(\sigma y)=1$.

Thus $\sigma^{-1}\{r=0\}\subseteq\{r=0\}$ and
$\{r=1\}\subseteq\sigma^{-1}\{r=1\}$ modulo null sets.
Probability preservation makes both inclusions equalities modulo null
sets.  Ergodicity makes each event null or conull.  The trace formula and
$0<\tau(p)<1$ exclude both conull cases.
\end{proof}

\subsection{A compact-set transfer lemma}

Projection fields need not be continuous.  The following observation
allows us to pass an identity to a weak limit of endpoint measures, provided
that their marginals do not concentrate on small sets.  A family of
finite measures is \emph{uniformly absolutely continuous} with respect to
$\mu$ if, for each $\varepsilon>0$, one can choose $\delta>0$ such that every
set of $\mu$-measure below $\delta$ has measure below $\varepsilon$ for every
member of the family.

\begin{lemma}\label{lem:compact-transfer}
Let $X$ be a compact metric space with Borel probability measure $\mu$.
Let $p(x)$ be a Borel field of orthogonal projections on $\ell^2(\mathbb Z)$,
meaning that every column $x\mapsto p(x)e_k$ is Borel, and let
$F_k:X\to(0,\infty)$, $k\in\mathbb Z$, be finite Borel functions.  Suppose
finite positive Borel measures $\beta_l$ on $X^2$ converge weakly to $\beta$.
Assume that both families of marginals of $\beta_l$ are uniformly absolutely
continuous with respect to $\mu$, and that both marginals of $\beta$ are
absolutely continuous.  Suppose also that some $0<B<\infty$ satisfies
\[
 F_k(y)/F_k(x)\le B\quad(k\in\mathbb Z)
\]
almost everywhere for each $\beta_l$ and for $\beta$.  If
\begin{equation}\label{eq:compact-transfer-identity}
 (1-p(y))\operatorname{diag}_k\bigl(F_k(y)/F_k(x)\bigr)p(x)=0
\end{equation}
holds $\beta_l$-almost everywhere for every $l$, then it holds
$\beta$-almost everywhere.
\end{lemma}

\begin{proof}
First replace every diagonal entry by its minimum with $B$, so that the
diagonal field is bounded everywhere.  Given an arbitrarily small
$\varepsilon>0$, simultaneous Lusin approximation supplies a compact
$Q\subset X$ on which every $F_k$ and every vector map $x\mapsto p(x)e_k$
is continuous, while $\mu(X\setminus Q)$ is sufficiently small to make
the $\beta_l$-mass outside $Q^2$ less than $\varepsilon$ for all $l$.
For completeness, the simultaneous assertion follows by applying Lusin's
property to the countably many maps with summable exceptional measures
and intersecting the resulting compact sets.  The usual property for
maps into separable metric spaces follows by finite-valued approximation
and measure regularity.

On $Q^2$, the capped diagonal field is strongly continuous: continuity on
the basis vectors and its uniform bound imply continuity on every vector.
The same argument applies to $p(x)$ and $p(y)$.  Products of strongly
continuous, uniformly bounded fields are strongly continuous.  Therefore
the zero set $Z_Q$ of the capped version of
\eqref{eq:compact-transfer-identity} in $Q^2$ is compact.  The assumed
identity implies
$\beta_l(Z_Q)\ge\beta_l(X^2)-\varepsilon$.  The compact-set upper bound
under weak convergence and convergence of total masses give
\[
 \beta(Z_Q)\ge\limsup_l\beta_l(Z_Q)
 \ge\beta(X^2)-\varepsilon.
\]
The capped global zero set thus has full $\beta$-measure, since
$\varepsilon$ was arbitrary.  It is measurable by its countably many
matrix entries.  Finally the cap changes no diagonal entry
$\beta$-almost everywhere.
\end{proof}

\subsection{The measurable weight and its cocycles}

We return to the prescribed rotation on $X=\mathbb R/\mathbb Z$.
The next formulas supply its positive weight and the diagonal ratios
to which Lemma~\ref{lem:compact-transfer} will apply.
For the parameters of Proposition~\ref{prop:parameters}, put
\begin{equation}\label{eq:measurable-weight}
 g_*(x)=\sum_j C_j\sum_{i\in J_j}
       \lfloor u_i(x)+\alpha_i\rfloor,\qquad
 f_*(x)=e^{-g_*(x)},\qquad T_*=U\pi(f_*).
\end{equation}
The sum defining $g_*$ has only finitely many nonzero terms almost
everywhere.  Here $c_i=C_j$ for $i\in J_j$ in the introduction's notation.
To prove the finiteness assertion, observe that
\[
 \sum_i\mu\{\lfloor u_i+\alpha_i\rfloor\ne0\}
 =\sum_i\alpha_i\le2\sum_j K_j^{-1}<\infty.
\]
The elementary first Borel--Cantelli implication follows here simply by
bounding the measure of a union of tail events by the sum of their measures.
Thus $0<f_*\le1$ almost everywhere.

For each $k\in\mathbb Z$, define
\begin{equation}\label{eq:cocycle-definition}
 F_k(x)=\exp\left(-\sum_j C_j\sum_{i\in J_j}
                     \lfloor u_i(x)+k\alpha_i\rfloor\right).
\end{equation}
These sums, too, have only finitely many nonzero terms almost everywhere:
in the tail where $|k|\alpha_i<1$, the probability of a nonzero floor is
$|k|\alpha_i$.  After taking a countable intersection, there is an invariant
conull set on which all $F_k$ are finite and strictly positive and
\begin{equation}\label{eq:cocycle-identities}
 F_{v+n}(x)=F_v(x)F_n(\sigma^v x),\qquad F_0=1,\qquad F_1=f_*.
\end{equation}
Indeed $u_i(\sigma^v x)=\{u_i(x)+v\alpha_i\}$, so the addition formula for
floors proves the identity term by term.  Since $\alpha_i>0$, $F_k(x)$ is
nonincreasing in $k$.  Give these functions arbitrary positive finite Borel
values on the remaining null set.

For endpoints in the common conull set, write
\[
 R_{x,y}(k)=F_k(y)/F_k(x).
\]
If $u_i(x)\le u_i(y)$, the $i$th coordinate contributes at most $1$ to
this ratio.  For $i\in J_s$, it contributes $e^{-C_s}$ precisely when
\begin{equation}\label{eq:ratio-crossing}
 u_i(x)<1-\{k\alpha_i\}\le u_i(y).
\end{equation}
When $k\alpha_i$ is integral the displayed event is empty, as it should be.
Without an order assumption, an individual coordinate contributes at most
$e^{C_s}$, since the two floors differ by at most one.

\subsection{From invariant projections to endpoint identities}

Suppose $p\in M$ is a projection satisfying $(1-p)T_*p=0$.  Its range is
then invariant under every positive power of $T_*$.  The formula
\[
 T_*(x)^n=S^n\operatorname{diag}_v F_n(\sigma^v x)
\]
and covariance \eqref{eq:field-covariance} show that, for $y=\sigma^n x$,
\begin{equation}\label{eq:positive-orbit-identity}
 (1-p(y))\operatorname{diag}_v R_{x,y}(v)p(x)=0\qquad(n\ge1).
\end{equation}
To verify the division involved, the cocycle identity gives
$F_n(\sigma^v x)/F_n(x)=F_v(\sigma^n x)/F_v(x)$; the positive divisor
$F_n(x)$ is finite.  The resulting diagonal is bounded on each such fiber
by $1/F_n(x)$, since $F_n\le1$.  This bound need not be uniform in $x$.

If $0<\tau(p)<1$, then
\begin{equation}\label{eq:diagonal-interior}
 0<r(x):=\langle p(x)e_0,e_0\rangle<1
 \quad\text{for almost every }x.
\end{equation}
This is Lemma~\ref{lem:positive-diagonal}, applied with $f=f_*$,
the faithful trace \eqref{eq:field-trace}, and covariance
\eqref{eq:field-covariance}.

The next lemma is the payoff from the sampling construction.  It extends
\eqref{eq:positive-orbit-identity} beyond actual orbit pairs, while retaining
a bound on every diagonal entry at once.

\begin{lemma}\label{lem:rational-graphs}
Let $p\in M$ be a projection satisfying $(1-p)T_*p=0$.
For every integer $A\ge1$ and every $t\in X$ that is an integer multiple of
$1/q_A$, the identity
\begin{equation}\label{eq:rational-graph-identity}
 (1-p(x+t))\operatorname{diag}_v R_{x,x+t}(v)p(x)=0
\end{equation}
holds for almost every $x$.  The exceptional sets may be chosen
simultaneously for this countable family of translations.
\end{lemma}

\begin{proof}
Fix $A,t$.  In every sufficiently late stage $j$, put $a=d_{j-1}\ge A$.
Since $q_A$ divides $q_a$, the sampling set $\mathcal N_{j,t}$ occurs in
Proposition~\ref{prop:parameters}.  Write $w=w_{m_j}$ and
$w^{(2)}=w_{m_j}^{(2)}$ for this stage's kernels in
\eqref{eq:sampling-kernel}.  Average over $n\in\mathcal N_{j,t}$
the graph measures $x\mapsto(x,\sigma^n x)$ with starting density
\[
 \prod_{i\in J_j}w\bigl(u_i(x),\{n\alpha_i\}\bigr).
\]
By \eqref{eq:sampling}, both marginal densities tend to $1$ in $L^2(\mu)$.
For the second marginal this uses the second orientation
$w^{(2)}(z,v)=w(z-v,v)$.  As every translation $n\theta$ lies within
$\ell_j$ of $t$ and $\ell_j\to0$, testing continuous functions on $X^2$
shows that these measures converge weakly to the Haar-first graph measure
of $x\mapsto x+t$.  At every point of positive density,
\[
 0<\{n\alpha_i\}<m_j^{-3/5}\Delta(u_i(x))\quad(i\in J_j),
\]
so there is no wrap and all current-block coordinates are ordered.

For every selected $n$, remove the same set of starting points, namely
those satisfying at least one of
\[
 \begin{array}{ll}
 u_i(x)\ge1-q_a\ell_j & (A\le i\le a),\\
 u_i(x)\ge1-2Q_j/N_s & (i\in J_s,\ s>j).
 \end{array}
\]
The removed Haar measure is at most
\[
 a q_a\ell_j+2Q_j\sum_{s>j}K_s^{-1}
 \le (j+1)^{-2}+4Q_j/K_{j+1}\longrightarrow0.
\]
The $L^2$ bound on the averaged first marginal and Cauchy--Schwarz show that the
removed weighted mass tends to zero.  Thus the trimmed measures have the
same weak limit.  Each trimmed marginal is dominated by its untrimmed
counterpart, so both retain uniform $L^2$ bounds and hence uniform absolute
continuity.

For $A\le i\le a$, the number $q_i t$ is integral.  The coordinate advance
is therefore $q_i\{n\theta-t\}<q_a\ell_j$, and trimming prevents a wrap.
For future coordinates it is $n\alpha_i<2Q_j/N_s$, again without a wrap.
Together with the current block, this orders every coordinate $i\ge A$.
Hence, on all trimmed measures,
\[
 R_{x,y}(v)\le
 \exp\left(\sum_{i<A}C_{s(i)}\right)\quad(v\in\mathbb Z),
 \qquad i\in J_{s(i)}.
\]
The bound is independent of $j$ and $v$.  It also holds on the limiting
graph, because $u_i(x+t)=u_i(x)$ for $i\ge A$.  The positive-orbit identity
holds on the approximants, so Lemma~\ref{lem:compact-transfer} gives
\eqref{eq:rational-graph-identity}.  There are only countably many pairs
$(A,t)$ under consideration.
\end{proof}
